\section{Conclusion and Future Work}

\subsection{Conclusion}

This thesis presented a citation-grounded AI assistant for Vietnamese labor-law question answering. The system is built around a scoped legal corpus and is designed to retrieve legal evidence before generating an answer. Its main purpose is to help users access labor-law information with traceable legal bases, rather than to provide unrestricted legal advice.

The implemented system combines several components into one complete workflow. It processes selected Vietnamese labor-law documents into structured evidence chunks, stores them in a Qdrant hybrid index, represents legal relationships in a Neo4j graph, retrieves and assembles relevant contexts, generates grounded answers, and validates citations before returning the final response. The system is also deployed as a web application with a FastAPI backend, a React/Vite frontend, Supabase authentication and data storage, and trace logging for evaluation.

The evaluation shows that graph-augmented retrieval improves the system's ability to recover relevant legal evidence. Recall@10 increases from 0.735 with hybrid retrieval only to 0.835 when graph expansion is enabled. The final end-to-end pass rate is 0.780. The citation validation pass rate and the out-of-corpus refusal pass rate both reach 1.000, showing that the system is strong at keeping citations inside retrieved evidence and refusing unsupported questions within the benchmark.

However, the results also show clear limitations. Most failures come from retrieval gaps rather than answer generation. The system still struggles with paraphrased user questions, table and appendix lookup, hard-negative citation cases, and some cross-document procedure questions. These limitations confirm that legal question answering depends not only on the language model, but also on the quality of legal evidence retrieval and selection.

Overall, this thesis shows that a reliable legal QA assistant should be treated as a full software pipeline. For Vietnamese labor-law questions, preserving legal structure, retrieving related provisions, controlling graph expansion, and validating citations are all important for producing answers that users can inspect and trust.

\subsection{Future Work}

Future work should first improve retrieval robustness. Better query rewriting could help the system understand informal user questions and map them to the correct legal concepts. Table-aware and appendix-aware retrieval should also be improved, especially for questions that require retirement-age tables, forms, lists, or structured legal information.

Second, graph expansion should be refined. The current results show that Neo4j helps recover missing evidence, but it can also introduce related yet distracting provisions. Future versions should improve graph filtering, query-type classification, and ranking rules so that graph expansion adds only legally useful evidence for each question.

Third, the legal corpus should be expanded and versioned. The current system is limited to selected Vietnamese labor-law documents. A larger corpus should include additional decrees, circulars, amendments, official guidance, and documents related to wages, insurance, administrative penalties, foreign workers, taxation, and trade unions. Versioning is also important because legal documents may be amended, replaced, or become outdated over time.

Fourth, the evaluation should include expert legal review. Deterministic benchmark checks are useful for reproducible software evaluation, but they cannot fully judge legal interpretation. Future evaluation should combine automatic citation checks with expert-reviewed answers, real anonymized user questions, and temporal-validity checks.

Finally, the user-facing application can be improved with clearer citation display, better explanation of insufficient-context cases, and more transparent evidence inspection. These improvements would make the assistant easier to use for non-specialist users while keeping the system bounded to verified legal sources.

\subsection{Final Remark}

\enlargethispage{2\baselineskip}
The main lesson of this thesis is that legal question answering is not only a text generation problem. A fluent answer is not enough if the legal basis cannot be verified. For a Vietnamese labor-law assistant, reliability depends on the whole pipeline: structured legal corpus preparation, hybrid retrieval, graph-based evidence expansion, careful context selection, and strict citation validation.
